\documentclass[12pt]{article}
\usepackage[spanish]{babel}
\usepackage[utf8]{inputenc}
\usepackage[T1]{fontenc}
\usepackage{amsmath, amssymb}
\usepackage[a4paper, margin=2.5cm]{geometry}

\begin{document}

\section*{Ejercicio de interpretación}

Considere el siguiente modelo estimado:

\[
\log(consumo_i)= 3.5 - 0.9\,\log(precio_i) + 0.04\,ingreso_i - 0.0005\,ingreso_i^2 + 0.03\,hogar_i + u_i
\]

donde:

\begin{itemize}
    \item $consumo_i$: cantidad de café consumida mensualmente por el hogar $i$ (en kilogramos).
    \item $precio_i$: precio del café que enfrenta el hogar $i$ (en pesos por kilogramo).
    \item $ingreso_i$: ingreso mensual del hogar $i$, medido en cientos de miles de pesos.
    \item $hogar_i$: número de personas que viven en el hogar $i$.
    \item $u_i$: término de error.
\end{itemize}

\subsection*{Preguntas}

\begin{enumerate}
    \item[a)] Calcule el retorno marginal del consumo de café dado el ingreso.
    \item[b)] Calcule el retorno marginal del consumo de café respecto al precio.
    \item[c)] Calcule el retorno marginal del consumo de café dado el número de personas en el hogar.
    \item[d)] Explique cómo cambia el efecto del ingreso sobre el consumo a medida que aumenta el ingreso.
    \item[e)] \textquestiondown En qué nivel de ingreso los retornos marginales del consumo empiezan a ser negativos?
\end{enumerate}

\subsection*{Respuestas}

\begin{enumerate}
    \item[a)] El retorno marginal del consumo respecto al ingreso se obtiene derivando $\log(consumo_i)$ respecto de $ingreso_i$:

    \[
    \frac{\partial \log(consumo_i)}{\partial ingreso_i}
    = 0.04 - 0.001\,ingreso_i
    \]

    Por lo tanto, el efecto marginal del ingreso sobre el logaritmo del consumo depende del nivel de ingreso.

    Como la variable dependiente está en logaritmo y $ingreso_i$ está en niveles, este efecto puede interpretarse aproximadamente como un cambio porcentual en el consumo. Así, un aumento de una unidad en $ingreso_i$ incrementa el consumo en aproximadamente

    \[
    100(0.04 - 0.001\,ingreso_i)\%
    \]

    ceteris paribus.

    \item[b)] El retorno marginal del consumo respecto al precio se obtiene derivando respecto de $precio_i$:

    \[
    \frac{\partial \log(consumo_i)}{\partial precio_i}
    = -0.9 \cdot \frac{1}{precio_i}
    = \frac{-0.9}{precio_i}
    \]

    Como tanto la variable dependiente como el precio están en logaritmos, el coeficiente $-0.9$ también se interpreta como elasticidad precio:

    \[
    \text{un aumento de 1\% en el precio del café reduce el consumo en aproximadamente } 0.9\%
    \]

    ceteris paribus.

    \item[c)] El retorno marginal del consumo respecto al número de personas en el hogar se obtiene derivando respecto de $hogar_i$:

    \[
    \frac{\partial \log(consumo_i)}{\partial hogar_i} = 0.03
    \]

    Como la variable dependiente está en logaritmo y $hogar_i$ está en niveles, esto implica que una persona adicional en el hogar aumenta el consumo de café en aproximadamente un $3\%$, ceteris paribus.

    \item[d)] El efecto del ingreso sobre el consumo está dado por:

    \[
    \frac{\partial \log(consumo_i)}{\partial ingreso_i}
    = 0.04 - 0.001\,ingreso_i
    \]

    Esto implica que el efecto marginal del ingreso es decreciente: al aumentar el ingreso, el impacto adicional del ingreso sobre el consumo de café se va reduciendo.

    Esto ocurre porque el coeficiente del término cuadrático, $-0.0005$, es negativo, lo que genera una relación cóncava entre ingreso y consumo.

    \item[e)] Los retornos marginales empiezan a ser negativos cuando el efecto marginal se hace igual a cero. Entonces resolvemos:

    \[
    0.04 - 0.001\,ingreso_i = 0
    \]

    \[
    0.001\,ingreso_i = 0.04
    \]

    \[
    ingreso_i = \frac{0.04}{0.001} = 40
    \]

    Por lo tanto, los retornos marginales del consumo respecto al ingreso se hacen cero cuando el ingreso es igual a 40.

    Como $ingreso_i$ está medido en cientos de miles de pesos, esto corresponde a:

    \[
    40 \times 100{,}000 = 4{,}000{,}000
    \]

    Es decir, los retornos marginales del consumo de café empiezan a ser negativos a partir de un ingreso mensual de 4.000.000 de pesos.
\end{enumerate}

\end{document}